% Generado por sistema/backend ReporteTex.java -- no editar a mano.
% Datos: dataset de prueba (respuestas sinteticas)
\paragraph{Grado de acuerdo con estar satisfecho con la formación.}

\textit{Paso 1. Contar con la muestra y determinar su tamaño.} Cada respuesta se codifica de 1 (totalmente en desacuerdo) a 4 (totalmente de acuerdo). El tamaño de la muestra es $n = 100$.

\begin{mitabla}{Cálculo de la media y la desviación estándar de la satisfacción con la formación (escala 1--4) \datosprueba}
\begin{tabular}{lccccc}
\toprule
Respuesta & $x_i$ & $f_i$ & $f_i\,x_i$ & $(x_i-\bar{x})^2$ & $f_i\,(x_i-\bar{x})^2$ \\
\midrule
Totalmente en desacuerdo & 1 & 8 & 8.00 & 4.8841 & 39.0728 \\
En desacuerdo & 2 & 9 & 18.00 & 1.4641 & 13.1769 \\
De acuerdo & 3 & 37 & 111.00 & 0.0441 & 1.6317 \\
Totalmente de acuerdo & 4 & 46 & 184.00 & 0.6241 & 28.7086 \\
\midrule
$\Sigma$ & & 100 & 321.00 & & 82.5900 \\
\bottomrule
\end{tabular}
\end{mitabla}

\textit{Paso 2. Calcular la media muestral.}
\[
\bar{x} = \frac{\sum f_i\,x_i}{n} = \frac{321.00}{100} = 3.2100\ \text{puntos}
\]

\textit{Paso 3. Calcular la desviación estándar muestral} (Coss Bu, ecuación 6.5).
\[
S = \sqrt{\frac{\sum f_i\,(x_i-\bar{x})^2}{n-1}} = \sqrt{\frac{82.5900}{99}} = 0.9134
\]

\textit{Paso 4. Determinar el estadístico teórico.} Con $\alpha = 0.05$ y $n-1 = 99$ grados de libertad, como los grados de libertad son más de 30 la distribución $t$ se aproxima con la normal, igual que en la prueba de los promedios: $t_{99,\,0.025} \approx z_{0.025} = 1.96$.

\textit{Paso 5. Calcular el margen de error.}
\[
\frac{S}{\sqrt{n}}\,t_{n-1,\,\alpha/2} = \frac{0.9134}{\sqrt{100}}\times 1.960 = 0.0913 \times 1.960 = 0.1790
\]

\textit{Paso 6. Construir el intervalo y concluir.}
\[
3.2100 - 0.1790 \;\le\; \mu \;\le\; 3.2100 + 0.1790 \quad\Longrightarrow\quad 3.03 \le \mu \le 3.39
\]

Con un 95\% de confianza, la media de la satisfacción con la formación (escala 1--4) de todos los titulados está entre 3.03 y 3.39 puntos. \datosprueba

\paragraph{Concordancia entre la formación y los requerimientos del mercado.}

\textit{Paso 1. Contar con la muestra y determinar su tamaño.} Misma codificación de 1 a 4. El tamaño de la muestra es $n = 100$.

\begin{mitabla}{Cálculo de la media y la desviación estándar de la concordancia formación--mercado (escala 1--4) \datosprueba}
\begin{tabular}{lccccc}
\toprule
Respuesta & $x_i$ & $f_i$ & $f_i\,x_i$ & $(x_i-\bar{x})^2$ & $f_i\,(x_i-\bar{x})^2$ \\
\midrule
Totalmente en desacuerdo & 1 & 10 & 10.00 & 3.3856 & 33.8560 \\
En desacuerdo & 2 & 22 & 44.00 & 0.7056 & 15.5232 \\
De acuerdo & 3 & 42 & 126.00 & 0.0256 & 1.0752 \\
Totalmente de acuerdo & 4 & 26 & 104.00 & 1.3456 & 34.9856 \\
\midrule
$\Sigma$ & & 100 & 284.00 & & 85.4400 \\
\bottomrule
\end{tabular}
\end{mitabla}

\textit{Paso 2. Calcular la media muestral.}
\[
\bar{x} = \frac{\sum f_i\,x_i}{n} = \frac{284.00}{100} = 2.8400\ \text{puntos}
\]

\textit{Paso 3. Calcular la desviación estándar muestral} (Coss Bu, ecuación 6.5).
\[
S = \sqrt{\frac{\sum f_i\,(x_i-\bar{x})^2}{n-1}} = \sqrt{\frac{85.4400}{99}} = 0.9290
\]

\textit{Paso 4. Determinar el estadístico teórico.} Con $\alpha = 0.05$ y $n-1 = 99$ grados de libertad, como los grados de libertad son más de 30 la distribución $t$ se aproxima con la normal, igual que en la prueba de los promedios: $t_{99,\,0.025} \approx z_{0.025} = 1.96$.

\textit{Paso 5. Calcular el margen de error.}
\[
\frac{S}{\sqrt{n}}\,t_{n-1,\,\alpha/2} = \frac{0.9290}{\sqrt{100}}\times 1.960 = 0.0929 \times 1.960 = 0.1821
\]

\textit{Paso 6. Construir el intervalo y concluir.}
\[
2.8400 - 0.1821 \;\le\; \mu \;\le\; 2.8400 + 0.1821 \quad\Longrightarrow\quad 2.66 \le \mu \le 3.02
\]

Con un 95\% de confianza, la media de la concordancia formación--mercado (escala 1--4) de todos los titulados está entre 2.66 y 3.02 puntos. \datosprueba

\paragraph{Pertinencia de las labores del puesto actual.}

\textit{Paso 1. Contar con la muestra y determinar su tamaño.} Solo responden quienes trabajan en una organización; misma codificación de 1 a 4. El tamaño de la muestra es $n = 68$.

\begin{mitabla}{Cálculo de la media y la desviación estándar de la pertinencia del puesto (escala 1--4) \datosprueba}
\begin{tabular}{lccccc}
\toprule
Respuesta & $x_i$ & $f_i$ & $f_i\,x_i$ & $(x_i-\bar{x})^2$ & $f_i\,(x_i-\bar{x})^2$ \\
\midrule
Totalmente en desacuerdo & 1 & 8 & 8.00 & 3.6548 & 29.2388 \\
En desacuerdo & 2 & 15 & 30.00 & 0.8313 & 12.4697 \\
De acuerdo & 3 & 20 & 60.00 & 0.0078 & 0.1557 \\
Totalmente de acuerdo & 4 & 25 & 100.00 & 1.1843 & 29.6064 \\
\midrule
$\Sigma$ & & 68 & 198.00 & & 71.4706 \\
\bottomrule
\end{tabular}
\end{mitabla}

\textit{Paso 2. Calcular la media muestral.}
\[
\bar{x} = \frac{\sum f_i\,x_i}{n} = \frac{198.00}{68} = 2.9118\ \text{puntos}
\]

\textit{Paso 3. Calcular la desviación estándar muestral} (Coss Bu, ecuación 6.5).
\[
S = \sqrt{\frac{\sum f_i\,(x_i-\bar{x})^2}{n-1}} = \sqrt{\frac{71.4706}{67}} = 1.0328
\]

\textit{Paso 4. Determinar el estadístico teórico.} Con $\alpha = 0.05$ y $n-1 = 67$ grados de libertad, como los grados de libertad son más de 30 la distribución $t$ se aproxima con la normal, igual que en la prueba de los promedios: $t_{67,\,0.025} \approx z_{0.025} = 1.96$.

\textit{Paso 5. Calcular el margen de error.}
\[
\frac{S}{\sqrt{n}}\,t_{n-1,\,\alpha/2} = \frac{1.0328}{\sqrt{68}}\times 1.960 = 0.1252 \times 1.960 = 0.2455
\]

\textit{Paso 6. Construir el intervalo y concluir.}
\[
2.9118 - 0.2455 \;\le\; \mu \;\le\; 2.9118 + 0.2455 \quad\Longrightarrow\quad 2.67 \le \mu \le 3.16
\]

Con un 95\% de confianza, la media de la pertinencia del puesto (escala 1--4) de todos los titulados está entre 2.67 y 3.16 puntos. \datosprueba

\paragraph{Aspectos de la Carrera, asignaturas.} Son preguntas de elección múltiple: cada titulado puede marcar varias opciones, por lo que $f_i/n$ se calcula sobre el número de titulados ($n = 100$) y los porcentajes pueden sumar más de 100.

\begin{mitabla}{Aspectos de mayor utilidad ($n=100$) \datosprueba}
\begin{tabular}{lccc}
\toprule
Aspecto & $f_i$ & $f_i / n$ & \% \\
\midrule
Distribución de horarios & 55 & $55/100 = 0.550$ & 55.0 \\
Equipos multidisciplinarios & 55 & $55/100 = 0.550$ & 55.0 \\
Pasantías o prácticas & 49 & $49/100 = 0.490$ & 49.0 \\
Habilidades analíticas & 47 & $47/100 = 0.470$ & 47.0 \\
Aplicación práctica en proyectos reales & 46 & $46/100 = 0.460$ & 46.0 \\
Herramientas de la industria & 43 & $43/100 = 0.430$ & 43.0 \\
\midrule
Menciones & 295 & & 295.0 \\
\bottomrule
\end{tabular}
\end{mitabla}

\textit{Lectura:} la respuesta más frecuente es ``Distribución de horarios'' ($f_i = 55$, 55.0\% de 100).

\begin{mitabla}{Aspectos que pueden mejorarse ($n=100$) \datosprueba}
\begin{tabular}{lccc}
\toprule
Aspecto & $f_i$ & $f_i / n$ & \% \\
\midrule
Distribución de horarios & 52 & $52/100 = 0.520$ & 52.0 \\
Actividades extracurriculares & 49 & $49/100 = 0.490$ & 49.0 \\
Relación con empresas (prácticas) & 48 & $48/100 = 0.480$ & 48.0 \\
Actualizar el plan de estudios & 48 & $48/100 = 0.480$ & 48.0 \\
Integrar teoría y práctica & 46 & $46/100 = 0.460$ & 46.0 \\
Más opciones de especialización & 43 & $43/100 = 0.430$ & 43.0 \\
\midrule
Menciones & 286 & & 286.0 \\
\bottomrule
\end{tabular}
\end{mitabla}

\textit{Lectura:} la respuesta más frecuente es ``Distribución de horarios'' ($f_i = 52$, 52.0\% de 100).

\begin{mitabla}{Asignaturas con ventaja competitiva ($n=100$) \datosprueba}
\begin{tabular}{lccc}
\toprule
Asignatura & $f_i$ & $f_i / n$ & \% \\
\midrule
Base de Datos & 65 & $65/100 = 0.650$ & 65.0 \\
Redes de Computadora & 60 & $60/100 = 0.600$ & 60.0 \\
T.I.S. & 59 & $59/100 = 0.590$ & 59.0 \\
Básicas & 57 & $57/100 = 0.570$ & 57.0 \\
Introducción & 55 & $55/100 = 0.550$ & 55.0 \\
Planif. de Proyectos y Gestión & 55 & $55/100 = 0.550$ & 55.0 \\
Tecnología Redes Avanzadas & 52 & $52/100 = 0.520$ & 52.0 \\
Inteligencia Artificial & 49 & $49/100 = 0.490$ & 49.0 \\
Electivas & 48 & $48/100 = 0.480$ & 48.0 \\
\midrule
Menciones & 500 & & 500.0 \\
\bottomrule
\end{tabular}
\end{mitabla}

\textit{Lectura:} la respuesta más frecuente es ``Base de Datos'' ($f_i = 65$, 65.0\% de 100).

\begin{mitabla}{Asignaturas poco útiles ($n=100$) \datosprueba}
\begin{tabular}{lccc}
\toprule
Asignatura & $f_i$ & $f_i / n$ & \% \\
\midrule
Contabilidad Básica & 58 & $58/100 = 0.580$ & 58.0 \\
Taller Prog. Bajo Nivel & 49 & $49/100 = 0.490$ & 49.0 \\
Básicas & 48 & $48/100 = 0.480$ & 48.0 \\
Aplic. Interactivas TV Digital & 47 & $47/100 = 0.470$ & 47.0 \\
Programación Funcional & 47 & $47/100 = 0.470$ & 47.0 \\
Teoría de Grafos & 47 & $47/100 = 0.470$ & 47.0 \\
Telefonía IP & 45 & $45/100 = 0.450$ & 45.0 \\
Teoría de Autómatas & 43 & $43/100 = 0.430$ & 43.0 \\
Web Semánticas & 40 & $40/100 = 0.400$ & 40.0 \\
Robótica & 39 & $39/100 = 0.390$ & 39.0 \\
Graficación por Computadora & 37 & $37/100 = 0.370$ & 37.0 \\
\midrule
Menciones & 500 & & 500.0 \\
\bottomrule
\end{tabular}
\end{mitabla}

\textit{Lectura:} la respuesta más frecuente es ``Contabilidad Básica'' ($f_i = 58$, 58.0\% de 100).

